\section{Open problems}\label{sec:open}

The question~(T) is open. For oblivious machines it is the question about the function $r(\eps)$ of
\cref{sec:one-function}. The problems below are the precise points at which this paper stops; the first three are finite
or algebraic and can be attacked directly.

\begin{problem}[Classes with no almost-period]\label{prob:aperiodic}
For a partition $\mathcal P$ of $\Z/d$ into $M$ classes let $D(\mathcal P)$ be the largest dimension of a code of a
class-uniform rule (\cref{def:class-uniform}), over a given field. \Cref{thm:TP} gives $D(\mathcal P)\le(d^2+Md)/2$ when
the classes are intervals. For general classes this fails (\cref{prop:aperiodic}), and for every fixed number $M\ge2$
of classes it fails by an amount that grows with $d$ (\cref{prop:two-class,cor:every-M}).
\begin{enumerate}[label=(\alph*),leftmargin=2em]
\item For two classes, is $D(\mathcal P)\le(\tfrac12+o(1))d^2$? More generally, is this so for $o(d)$ classes?
For two classes of any shape the best upper bound known is
$\tfrac{41}{64}d^2+\tfrac53d^{3/2}$~\cite[Theorem~3.4]{SharedDecoders2026}; before that paper it was
$\tfrac34d^2+O(d)$, from \cref{thm:FL}, together with what \cref{thm:TP} gives for the best shift set. For $M\le d/2$
classes that paper proves the bound
$\bigl(1-1/(4\lceil\sqrt{2M}\,\rceil)\bigr)d^2$~\cite[Theorem~3.5]{SharedDecoders2026}, which is below the value
$(1-\tfrac1{2M})d^2+O(d)$ of \cref{thm:FL} when $M\ge11$. The largest gains known for a fixed number of
classes are of order~$d$.
\item Determine $D$ for the classes $\{0,1,3\},\{2,4,5\}$ of $\Z/6$. By \cref{thm:TP} it is at most $28$ over every
field. It is at least $25$ over every field in which $3\ne0$: over $\F_2$ by \cref{prop:aperiodic}, and in general
by~\cite[Theorem~3.2]{ClassUniformMemo2026}. In characteristic~$3$ no code of dimension $25$ is known.
\item Fix the number $M\ge2$ of classes. Is there a constant $C_M$ such that $D(\mathcal P)\le(d^2+Md)/2+C_Md$ for
every $d$ and every partition $\mathcal P$ of $\Z/d$ into $M$ classes? Or is there a family of partitions into $M$
classes, with codes that have a description that can be analysed, in which $D(\mathcal P)-(d^2+Md)/2$ grows faster
than~$d$?
\end{enumerate}
\end{problem}

A positive answer to (a) would be the half bound for all piecewise-uniform decoders in the one-hop model; it needs a
count sharper than the union bound of \cref{thm:TP}, which loses a factor $s/(s+1)$ for $s+1$ translates. A negative
answer would be a coding gain of positive rate over the pairing value from $o(d)$ classes, native to masks and
rotations, the first of its kind, and its cost on a tape would become the question. \Cref{sec:small} records what exact
computation says for small~$d$.

The two bounds quoted in~(a) come from a later paper of the series~\cite{SharedDecoders2026}, which proves them in
Lean, over every field, for codes that are subspaces and decoders that are arbitrary maps (\cref{sec:sets}). Its count
is sharper than the union bound in the way asked for above. The code is intersected with one translate, the
intersection lies in the space of matrices that are covariant on the components of that shift, this space is
invariant under rotation, and the next intersection is taken inside it; so with shift sets $E_0=\{0\}$ and
$E_j=E_{j-1}\cup(E_{j-1}+t_j)$ for $1\le j\le k$, a code has dimension at most
$\sum_{j<k}n_j/2^{j+1}+n_k/2^k$, where $n_j$ is $d$ times the number of components for
$E_j$~\cite[Theorem~3.2]{SharedDecoders2026}. When $E_2$ has four elements, \cref{thm:TP} for $E_2$ gives
$\tfrac34n_0+\tfrac14n_2$ where this gives $\tfrac12n_0+\tfrac14n_1+\tfrac14n_2$, and $n_0=d^2$. The value $\tfrac{41}{64}$ comes from the shifts $s,3s,9s$ and an average over~$s$. For
almost every partition into two classes the same paper has a slightly smaller constant, and it proves that the
certificates of this method cannot go below $(\tfrac58-o(1))d^2$, and those of a variant not below
$(\tfrac9{16}-o(1))d^2$, for almost every such partition~\cite[Theorems~C.5, C.6 and~C.7]{SharedDecoders2026}. So (a) is open, and going below
these values needs more than the certificates of that method.

Part~(c) has a history. Revision \texttt{b2c389a} of this paper asked only for a family of partitions and codes, with
a description that can be analysed, in which $D(\mathcal P)-(d^2+Md)/2$ grows, without fixing the number of classes.
The memo~\cite{ClassUniformMemo2026} answered that question as it was worded: the blow-ups of the code of
\cref{prop:aperiodic} are codes of dimension $25k^2$ for $2k$ classes on $\Z/6k$, so that
$D(\mathcal P)-(d^2+Md)/2\ge k^2$ there~\cite[Theorem~4.2 and Corollary~4.4]{ClassUniformMemo2026}. The family is
described given its base code, which a solver found. A blow-up multiplies the excess of a code over $(d^2+Md)/2$ and
the product $dM$ by the same factor~\cite[Corollary~4.3]{ClassUniformMemo2026}, so the blow-ups of one finite example
gain only together with a growing number of classes. The memo and revision \texttt{14b157a} of this paper therefore
asked whether $D(\mathcal P)-(d^2+Md)/2$ is bounded for a fixed number of
classes~\cite[Problem~6.2]{ClassUniformMemo2026}. It is not. An update of the memo showed this for $M=2$ over every
field in which $2\ne0$, with an explicit rule whose code has a gain of
$d/4-1$~\cite[Theorem~4.8 and Corollary~4.10]{ClassUniformMemo2026}, and left open characteristic~$2$ and every
$M\ge3$. \Cref{prop:two-class} is that theorem with its two coefficients $2$ and $\tfrac12$ replaced by $\lambda$ and
$\lambda^{-1}$, which makes it hold over every field. \Cref{cor:every-M} blows the two-class codes up by a fixed
factor, with a code of its own at each position, and gives an unbounded gain for every $M\ge2$ over every field. What
remains is the rate of growth, which is~(c) as now stated; for $M=2$ it is, in substance,
\cite[Problem~6.3]{ClassUniformMemo2026}.

By \cref{cor:every-M} a constant $C_M$ as in~(c) is at least $\lfloor M/2\rfloor/4$, over every field; for two classes
it is at least $\tfrac14$. For the two classes of \cref{prop:two-class}, \cref{thm:TP} bounds the gain by $d/2$. In
general \cref{thm:TP} does not answer~(c): with two shifts it bounds the gain of a two-class partition by
$(\chi-2)d/2$, where the number $\chi$ of components can be of order~$d$, and with three or more shifts its bound on
$D(\mathcal P)$ is at least $2d^2/3$~\cite[Section~6]{ClassUniformMemo2026}. A positive answer to~(c) for $M=2$ would
prove (a) for two classes. A family of two-class partitions with a gain of order $d^2$ would be a negative answer
to~(a).

For partitions with a bounded period the answer to~(c) is positive, and for some of them the constant is known. Let a
partition $p$ of $\Z/s$ be lifted to the partition $p^\uparrow$ of $\Z/d$, $d=ms$, in which $c$ has the class of
$c\bmod s$. \Cref{thm:TP} with the shifts $0$ and $s$ bounds its largest code by $(d^2+sd)/2$. If $m$ is prime to $s$
and not divisible by the characteristic, and the field contains the $m$-th roots of unity, the exact value is
\[
  D(p^\uparrow)=\frac{d^2-ms^2}2+m\,D(p)=\frac{d^2}2+\frac ds\Bigl(D(p)-\frac{s^2}2\Bigr);
\]
over a field without these roots of unity the right side is a lower bound, and it is an upper bound when $D(p)$ is
taken over the algebraic closure~\cite[Theorems~5.1 and~5.3]{SharedDecoders2026}. So the gain of a partition with
period $s$ is at most $(s-M)d/2$, for these lifts it is a fixed multiple of~$d$, which is zero when $D(p)$ is the
value $(s^2+Ms)/2$ of intervals, and a gain that grows faster than $d$ needs partitions without a bounded period; the
classes of \cref{prop:two-class} have no period smaller than~$d$. From the classes of \cref{prop:aperiodic} the lift
gives two classes on $\Z/6m$, for $m$ prime to~$6$, with a binary code of dimension $d^2/2+7d/6$, and this is the
whole code of the lifted rule~\cite[Corollary~5.2 and Theorem~B.2]{SharedDecoders2026}. Its gain over $(d^2+2d)/2$
is $d/6$; the rule of \cref{prop:two-class} gains $d/4-1$, at the sizes divisible by four.

\begin{problem}[Unequal pieces]\label{prob:unequal}
For which partitions of $\Z/d$ into $M$ intervals, and over which fields, is the bound $\lfloor(d^2+Md)/2\rfloor$ of
\cref{thm:TP} attained? It is attained for equal pieces (\cref{prop:block-symmetric}). Over $\F_2$ it is attained in
every case computed in \cref{sec:small} except one, the pieces $\{0\},\{1\},\{2,3,4\}$ at $d=5$, where the largest code
has dimension $19<20$ on the word of a solver. Some of the codes that attain it are not codes of copy rules.
\end{problem}

The question is about linear decoders. If the decoders may be arbitrary maps and the code is still a subspace, the
bound is attained over $\F_2$ for every partition into $M$ intervals with $d-M$ even, and the same construction gives
a code of dimension $(d^2+Md)/2$ for every partition into $M$ classes of any shape with $d-M$
even~\cite[Theorem~4.3]{SharedDecoders2026}. For the pieces $\{0\},\{1\},\{2,3,4\}$ at $d=5$ this is a code of
dimension $20$ with decoders that are not linear, where linear decoders reach $19$ on the word of a solver.

\begin{problem}[Both sides, all scales]\label{prob:both}
Let both the $A_a$ and the $B_a$ lie in the span of masks of dyadic intervals composed with rotations, with at most $M_j$
intervals of length $d/2^j$ in use at scale~$j$. Bound $\dim\ker(\sw-K)$ in terms of the profile $(M_j)_j$. The folklore
recursion shows that the full profile suffices for dimension $d^2$; \cref{thm:TP} and \cref{cor:digit} treat one side
and one scale at a time.
\end{problem}

\begin{problem}[Shared decoders on cone systems]\label{prob:machines}
\begin{enumerate}[label=(\alph*),leftmargin=2em]
\item State and prove a version of \cref{thm:TP} for the cone systems of oblivious linear-time schedules: on a stretch
of periodic head motion the outputs share their decoder, and a revealed bit or an irregular step changes the decoders
of some outputs. Does a bound of the form $r(\eps)\le\tfrac12+C(\eps+\eta)$ hold for schedules whose motion is
periodic outside a set of steps of density~$\eta$? (\Cref{sec:to-machines}.)
\item Fix a window $w\ge3$ and an alphabet of $q$ symbols. Does every storage code on a one-class relation of
window~$w$ (\cref{sec:sets}) have at most $q^{(1/2+o(1))d^2}$ elements as $d$ grows, or even at most $q^{d^2/2+Cd}$
with a constant $C$ that depends on $w$ and~$q$? And is this so at least when one function serves all rows, that is,
when the function for the position $(c,r)$ does not depend on~$c$? (For $w\le2$ the answer to both is yes, with
$C=\tfrac32$: \cref{prop:order}(b). For $w=2$ and a function for each position the least constant lies between $1$
and $\tfrac32$, by the examples after that proposition; for $w=1$ it is $\tfrac12$.)
\item Fix $w\ge3$ and an alphabet of $q$ symbols, and let $\rho$ assign to every column index at most $w$ rows. Does
every $\rho$-code (\cref{sec:sets}), that is, every set of matrices in which each column is a function of its $w$
rows, whole, have at most $q^{(1/2+o(1))d^2}$ elements as $d$ grows? (For $w\le2$ it has at most
$q^{d(d+1)/2+d-1}$, by the proof of \cref{prop:order}(b); in general at most $q^{wd(d+1)/(w+1)}$, by
\cref{prop:rows}(d).)
\end{enumerate}
\end{problem}

Part~(a) is the problem as the earlier revisions stated it; part~(b) was added in revision \texttt{3911d87}, for
every window, and is restricted here to the windows that remain open. The form of the bound
in~(a) comes from the one-hop model with affine rules. By \cref{cor:entry-advice}, over a finite field, if every rule
that the advice selects changes, from one row to the next, at no more than $E$ entries, then the matrices of a code
that share a value of the advice number at most $\abs F^{(d^2+d+E)/2}$, whatever the advice is. The rate of such a
set, the analogue of $r(\eps)$, is therefore at most $\tfrac12+(d+E)/(2d^2)$, and $E=O((\eps+\eta)d^2)$ gives a bound
of the form above. So the question is whether the revealed bits and the irregular steps of a schedule amount to that
many changes of decoders at single outputs, on the cone systems of real schedules and for decoders that are
compositions of a transition function. Revision \texttt{b2c389a} of this paper stated the count as ``a genie that
reveals $\eps L$ bits changes the regime at most $\eps L$ times'', and its \cref{sec:to-machines} priced a change of
regime at half a row. Half a row is the price of a change that affects all $d$ entries of a row, and at that price the
count gives nothing.

For decoders that are not affine the count of \cref{cor:entry-advice} is not available at all, and this is the
reason for part~(b). With a view of the whole matrix, one decoder for all rows, with no change of any kind, is correct
on a set of at least $q^{d^2}/(2d)$ matrices~\cite[Theorem~A.1 and Corollary~A.2]{SharedDecoders2026}
(\cref{sec:sets}). The inputs that share a value of advice form a set and not a subspace. So for them, and for
decoders that are compositions of a transition function, the sharing of a decoder bounds nothing when the decoder
sees the whole stored word. What that example uses, and the decoders of a machine do not have, is the size of the
view. The cones that a genie leaves have bounded size, so a decoder reads boundedly many bits, and with them it can
tell apart only boundedly many positions (\cref{sec:sets}). Whether the sharing of a decoder bounds a set of matrices
when the decoders read boundedly many symbols of a fixed alphabet is not known, and part~(b) asks it in the one-hop
model, in two forms. In the weaker form one function serves all rows and reads translated positions, which is what a
stretch of periodic motion gives. Since a set of matrices can contain marker rows, one function can then imitate
boundedly many classes of rows, where the marker rows are free: with markers the value $(d^2+d)/2$ of subspaces is
exceeded by a multiple of $d$ that grows with the window, over every finite field (\cref{sec:sets}), so the constant
$C$ of~(b) has to grow with~$w$. In the stronger form only
the positions are translates, and every position may have its own function. Over a fixed alphabet there are boundedly many functions of $w$ symbols, so the stronger form
is the question for boundedly many classes of any shape, drawn entry by entry. It contains what is left of the weaker
form among the inputs that share a revealed word: there the outputs of a periodic stretch whose cones meet the same
pattern of revealed bits, with the same values, compute the same function of the input bits that they read, so that
the revealed word divides the outputs into boundedly many classes of a shape that the advice chooses. The hypothesis
of the stronger form is a property of the relation alone, and so of the head motion. What is proved is in
\cref{prop:sets,prop:order,prop:mutual}: windows one and two; every window when the offsets, or the differences of
the rows from a permutation, are bounded or take boundedly many values; codes that use only helpers in both
directions; and, for the weaker form, subspaces. What these
results leave is a window of at least three on a relation whose rows and offsets are both spread. For small windows
the acyclic-subgraph bound, which gives all of them but the last, may still be enough; for large windows the computations
reported in \cref{sec:sets}, in the paragraph on what the order argument does not reach, indicate that it is not. A code above the pairing rate would have to use
helpers in one direction, on a relation of that kind.

The conclusion asked for in~(b) is that of the all-code question of the fifth
paper~\cite[Question~9.1]{DirectedRecovery2026}, with other hypotheses. That question is about relations of bounded
incidence degree that are separated, and for those the answer is no (\cref{sec:false}). A one-class relation has
the symmetry of its positions instead; it may have loops and cycles of length two, and one position may help $d$
others. We have not determined whether relations like those of the counterexamples of \cref{sec:false} can be
one-class relations; the undirected ones cannot, in the sense of \cref{prop:mutual}(b). In every one-class
relation of bounded window few positions are their own helpers, and $o(d^2)$ positions meet all transitive triangles
(\cref{prop:mutual}(c)). A rule with $M$ classes of any shape whose decoders have window $w$ gives a relation that is
contained in a one-class relation of window $Mw$. So a positive answer to the stronger form of~(b) would give the
half bound of \cref{prob:aperiodic}(a) for every fixed number of classes and decoders of bounded window over a fixed
finite field, for sets as well as subspaces. For subspaces the weaker form is enough. A code of a rule with $M$
classes whose decoders have window $w$ can be replaced by a code of at least the same dimension in which one row of
each class is free, with decoders of window at most
$2w+1$~\cite[Theorem~4.1]{SharedDecoders2026}; fixing $k=\lceil\log_qM\rceil$ of these rows as markers
(\cref{sec:sets}) gives a storage code with one function for all rows, on a one-class relation of window
$M(2w+1)+k$, with $q^{-kd}$ times as many elements. So a positive answer to the weaker form, for every window, gives
\cref{prob:aperiodic}(a) for rules with a fixed number of classes and decoders of bounded window over a fixed finite
field. A negative answer to either form would be a code above the pairing rate, with bounded windows, on positions
that move with the output row.

Part~(c) was added with \cref{prop:rows,prop:stages,prop:expanding,prop:pieces}. A storage code on a one-class
relation of window $w$ is a $\rho$-code with at most $w$ rows for each column, whatever its offsets, so a positive
answer to~(c) for every $w$ gives the bound $q^{(1/2+o(1))d^2}$ of~(b) for every window; it does not give the form
$d^2/2+Cd$. In the other direction (c) allows more than (b): an entry may depend on all entries of its rows. What is
proved is the case in which the rows of each column lie in a small block, and with it pipelines of stages that read
rows inside address cosets, where $t$ moved digits give exactly $(d^2+2^td)/2$; pipelines with buffered scans and
with differently arranged copies, more weakly (\cref{prop:stages}); and products of codes on boundedly many positions
(\cref{prop:pieces}). What is left are the relations on which the acyclic-subgraph bound stays a constant above one
half whatever order is used, among them those that expand and contain a bijection. There are such relations that are
unions of the identity and three increasing maps, so that an operator with these rows takes boundedly many forward
passes, given the strings that mark the maps (\cref{prop:expanding} and the paragraph after it). A positive answer
to~(c) for $w\ge4$ therefore needs an argument that is
not an order argument, and a negative answer on increasing maps, with functions computable in linear time, would be
an operator of boundedly many passes that transposes more than the pairing rate. By \cref{prop:pieces} a code above
the pairing rate cannot be a product of local codes. One family in which to look for either answer: for a group
$\Gamma$ of order $d$ and a bounded set $\Sigma\subseteq\Gamma$, index the rows and columns by $\Gamma$ and let the
helpers of $(c,r)$ be the positions $(rs,ct)$ for boundedly many pairs $(s,t)$ with $s\in\Sigma$; these relations
are contained in the relation of $\rho(r)=r\Sigma$, and for $\Gamma=\Z/d$ they are the one-class relations with
constant offsets and constant row differences, which \cref{prop:order}(d) settles.

For adaptive machines a later paper of the series puts the function $r(\eps)$ on the machine's own runs. At the
canonical entry width, for all but an exponentially small share of the inputs that a linear-time transposer handles,
and for every set $B$ of steps of the run on such an input, the relation between entries that is left when $B$ is cut
out carries storage codes whose deficit is at most a small multiple of the input length plus a constant times
$\card B$~\cite[Theorem~H.5 and Corollary~H.6]{AdaptiveTransfer2026}. So a bound on storage codes that is stated for
relations, as part~(b) is, can be applied to a single run of an adaptive machine wherever the relation of that run
satisfies its hypothesis.

\begin{problem}[What an unshuffle does not supply]\label{prob:scales}
\Cref{cor:digit} requires a transposer to break the symmetry of every address digit, and \cref{rem:compose-TP} shows
that commutator rank cannot turn this into a factor $\log d$, because one unshuffle breaks all digits at once. Find a
quantity that is subadditive over the steps of a machine, is large for the swap at every scale separately, and is not
raised at all scales by a single linear-time operation. This is the operator form of the multi-scale problem of the
sixteenth paper~\cite[Problem~11.1]{SixteenthPaper2026}.
\end{problem}

For pipelines whose stages read rows inside cosets of the address space, the dimension of the coset adds up over
the stages and bounds the set that is transposed, for all codes: one for each digit moved, and $\log_2d$ for the
swap (\cref{prop:stages}(a)). This is not the quantity asked for. It does not look at scales separately, and a
single linear-time operation can raise it to its maximum: a scan with a buffer of one row, or a stage that reads a
row $r$ together with its shuffle parents $\lfloor r/2\rfloor$ and $\lfloor r/2\rfloor+d/2$ (the remarks after
\cref{prop:stages}). For the
first the proposition has a price in advice, and for the second only the weaker count by the number of ones of the
address.

\begin{problem}[Relative rigidity on ball patterns]\label{prob:rigidity}
The eighth paper asks whether there is $\rho>0$ such that, for every linear-time schedule, every fixed $D$ and all
large~$d$, the transposition matrix is not the sum of a matrix of rank below $\rho L$ and a matrix supported on the
depth-$D$ ball pattern of the schedule~\cite[Problem~10.6]{EighthPaper2026}; a positive answer alone would exclude
oblivious linear-time transposers with linear data paths~\cite[Corollary~9.3]{BoundedStack2026}. The classes of this
paper are defined by structure on one tensor side. Which ball patterns of schedules are contained in such a class,
after a change of bases on one side? For those, \cref{thm:FL,thm:TP} give $\rho=\tfrac12-o(1)$.
\end{problem}
